\par\noindent\textbf{(1)} 变分给 \(-2A\theta''+2K\sin\theta\cos\theta=0\)，即 \(A\theta''=(K/2)\sin2\theta\)。乘 \(\theta'\) 并用无穷远边界 \(\theta'\to0,\sin\theta\to0\)，得到一次积分 \(A(\theta')^2=K\sin^2\theta\)。
\par\noindent\textbf{(2)} 选从0增至 \(\pi\) 的支 \(\theta'=\sqrt{K/A}\sin\theta\)，分离积分 \(\int d\theta/\sin\theta=\ln\tan(\theta/2)=(x-x_0)/\delta\)，其中 \(\delta=\sqrt{A/K}\)。因此 \(\theta=2\arctan e^{(x-x_0)/\delta}\)、\(\cos\theta=-\tanh[(x-x_0)/\delta]\)。
\par\noindent\textbf{(3)} 利用 \(A\theta'^2=K\sin^2\theta\)，能量 \(\sigma=2K\int_{-\infty}^\infty\sin^2\theta\,dx\)。因 \(\sin\theta=\operatorname{sech}((x-x_0)/\delta)\)，积分为 \(2\delta\)，得 \(\sigma=4K\delta=4\sqrt{AK}\)。交换能和各向异性能各为 \(2\sqrt{AK}\)。
\par\noindent\textbf{(4)} 交换刚度增大使 \(\delta\propto\sqrt A\) 变宽、壁能 \(\sigma\propto\sqrt A\) 增大；各向异性增大则壁变窄而总壁能仍增加。若加入Dzyaloshinskii–Moriya相互作用，一次积分与旋转手性会改变，不能直接沿用此无手性壁。
